\section{Teorema de confluencia de las respuestas dimensionales}
\label{rm:ch:confluencia-general}

\subsection{Antecedente com\'un y producto fibrado}

Los resultados del volumen proceden de lectores heterog\'eneos aplicados a
un mismo estado geneal\'ogico enriquecido.  Sea

\begin{equation}
 X=\bigl(X_{\APP},X_{\TRIT},X_{\TPK},
          \mathcal C,\mathcal B,\mathcal R,\mathcal M,\partial X\bigr)
 \label{rm:eq:confluence-state}
\end{equation}

el estado que conserva cilindros, acarreo, ruta, memoria y frontera. Sobre
\(X\) se construyen los lectores tipados

\begin{equation}
 P_i:X\longrightarrow Y_i,
 \qquad
 i\in I:=\{\alpha,{\rm act},{\rm ang},108,\square,{\rm obs},{\rm Per}\}.
 \label{rm:eq:confluence-objects}
\end{equation}

Sus codominios representan, respectivamente, el generador dodecaf\'asico de
acoplamiento, la pareja de acci\'on, el operador angular bicapa, el reloj
\(108\), la pantalla c\'ubica, el observador interno y el modo de Perron.
Para cada lector se toma su grafo, provisto de la proyecci\'on can\'onica

\begin{equation}
 \Gamma_i:=\{(x,y_i)\in X\times Y_i:P_i(x)=y_i\},
 \qquad
 \pi_i:\Gamma_i\longrightarrow X,
 \quad (x,y_i)\longmapsto x.
 \label{rm:eq:confluence-reader-graphs}
\end{equation}

El objeto maestro es el producto fibrado bien tipado de esos grafos:

\begin{equation}
 \boxed{
 \mathfrak M_X=
 \mathop{\times}_{X,\,i\in I}\Gamma_i
 =\left\{\bigl(x,(y_i)_{i\in I}\bigr):
 P_i(x)=y_i\ \text{para todo }i\in I\right\}.}
 \label{rm:eq:master-fiber-product}
\end{equation}

Las proyecciones \(\pi_i\) fijan materialmente el mismo antecedente. Un
elemento de \(\mathfrak M_X\) es una familia de salidas compatibles y
acompa\~nadas por el estado que las genera.

\begin{definition}[Publicaci\'on dimensional]
Una publicaci\'on dimensional es una composici\'on
\begin{equation}
 \mathfrak M_X\xrightarrow{\;P\;}\mathcal L
 \xrightarrow{\;\operatorname{ev}_{\mathcal U}\;}\RR,
 \label{rm:eq:dimensional-publication}
\end{equation}
donde \(\mathcal L\) es una recta dimensional tipada y
\(\operatorname{ev}_{\mathcal U}\) es la evaluaci\'on final en una carta de
unidades. El mapa \(P\) construye primero el invariante y su genealog\'ia.
\end{definition}

\subsection{Independencia horizontal y compatibilidad vertical}

Dos lectores hermanos comparten antecedente y pueden satisfacer relaciones
cruzadas sin factorizar uno por el otro.  Esta distinci\'on gobierna las
dobles hojas del vac\'io, las realizaciones positiva y orientada, la doble
proyecci\'on arquimediana y excepcional, y las lecturas geom\'etrica,
informacional y de sabor del tr\'ansito \(6\to8\).

\begin{theorem}[Confluencia sin factorizaci\'on horizontal]
Sean \(P_i:\mathfrak M_X\to Y_i\), \(i=1,\ldots,r\), publicaciones
construidas sobre el mismo estado. Sup\'ongase que cada \(P_i\) conserva sus
datos tipados y que las relaciones \(R_j(P_1,\ldots,P_r)=0\) se demuestran
despu\'es de construir los lectores.  Entonces el morfismo conjunto
\begin{equation}
 P=(P_1,\ldots,P_r):\mathfrak M_X\longrightarrow
 Y_1\times\cdots\times Y_r
 \label{rm:eq:joint-publication}
\end{equation}
es compatible verticalmente con \(X\). Una factorizaci\'on
\(P_i=F_{ij}\circ P_j\) requiere, adem\'as, que \(P_j\) separe todas las
fibras que \(P_i\) distingue. La mera existencia de \(R_j=0\) no implica
esa separaci\'on.
\end{theorem}

\begin{proof}
La compatibilidad vertical se obtiene por la propiedad universal del
producto fibrado de los grafos \(\Gamma_i\): todas las componentes poseen
la misma imagen en \(X\).
Si \(P_i=F_{ij}\circ P_j\), entonces
\(P_j(x)=P_j(x')\) fuerza \(P_i(x)=P_i(x')\). Por tanto, cualquier par
\(x,x'\) que comparta la lectura \(j\) y difiera en la lectura \(i\) refuta
la factorizaci\'on. Una relaci\'on escalar \(R_j=0\) s\'olo restringe la imagen
conjunta; no garantiza la inclusi\'on de n\'ucleos necesaria para factorizar.
\end{proof}

El teorema explica por qu\'e el producto de hojas puede fijar el cono causal
mientras el cociente conserva orientaci\'on, y por qu\'e el \`area total puede
coexistir con un Hamiltoniano modular que distingue la procedencia \(90/120\).

\subsection{1.ª confluencia: acoplamiento, vacío y sector electrodébil}

Sea

\begin{equation}
 D=D_A=\det T,
 \qquad
 Q(D)=-\frac9{25}\log D=4\pi\alpha_{\HMT}.
 \label{rm:eq:confluence-q-d}
\end{equation}

La memoria orientada se expresa mediante

\begin{equation}
 \rho=e^{2C^*_{\rm rad}},
 \qquad
 q_-=\sqrt{D\rho},
 \qquad
 q_+=\sqrt{D/\rho}.
 \label{rm:eq:confluence-qpm}
\end{equation}

El c\'alculo funcional

\begin{equation}
 F(z)=\frac{1-z^{90}}{1-z^{120}},
 \qquad
 r_\pm=F(q_\pm)
 \label{rm:eq:confluence-F}
\end{equation}

produce

\begin{equation}
 \widehat\mu=r_-^2,
 \quad
 \widehat\varepsilon=r_+^2,
 \quad
 \widehat Z=\frac{r_-}{r_+},
 \quad
 \widehat c=(r_-r_+)^{-1}.
 \label{rm:eq:confluence-vacuum}
\end{equation}

En la interfaz electrod\'ebil de \`arbol,

\begin{equation}
 g^2=\frac{Q(D)}{1-D},
 \qquad
 g'^2=\frac{Q(D)}{D},
 \label{rm:eq:confluence-gauge-couplings}
\end{equation}

de donde

\begin{equation}
 \boxed{
 (1-D)g^2=Dg'^2=Q(D),
 \qquad
 \frac1{g^2}+\frac1{g'^2}=\frac1{Q(D)}.}
 \label{rm:eq:confluence-ew-master}
\end{equation}

La primera igualdad publica el mismo acoplamiento en dos orientaciones de
calibre; la segunda lo recupera como suma de compliancias. El determinante
\(D\) gobierna la mezcla par y \(\rho\) conserva la orientaci\'on impar.

La coordenada de Higgs incorpora esta \`ultima:

\begin{equation}
 \lambda_H=\frac{Q(D)}{8(1-D)}D^{-3}\rho^{5/3}.
 \label{rm:eq:confluence-higgs}
\end{equation}

Definiendo

\begin{equation}
 \Xi=\frac{8(1-D)D^3\lambda_H}{Q(D)},
 \label{rm:eq:confluence-xi}
\end{equation}

se reconstruye

\begin{equation}
 \rho=\Xi^{3/5},
 \qquad
 \widehat Z=
 \frac{F(\sqrt D\,\Xi^{3/10})}
      {F(\sqrt D\,\Xi^{-3/10})}.
 \label{rm:eq:confluence-vacuum-higgs}
\end{equation}

La ecuaci\'on \cref{rm:eq:confluence-vacuum-higgs} constituye un control cruzado
entre orientaci\'on escalar e impedancia del vac\'io.  Su fuerza procede de la
construcci\'on paralela de ambos lectores sobre \(X\).

\subsection{2.ª confluencia: acción, gravedad y forma electrogravitatoria}

Sea

\begin{equation}
 Q_P=E_P,
 \qquad
 \ell_P,
 \qquad
 Q_P\ell_P=\hbar c,
 \qquad
 \frac{\ell_P}{Q_P}=\frac{G}{c^4}.
 \label{rm:eq:confluence-planck-pair}
\end{equation}

El producto es el cuanto de intercambio \(\hbar c\); el cociente es la
compliancia gravitatoria.  Introducimos

\begin{equation}
 \mathscr T_P=\frac{Q_P}{\ell_P}=\frac{c^4}{G},
 \qquad
 \mathscr C_P=\frac{\ell_P}{Q_P}=\frac{G}{c^4}.
 \label{rm:eq:confluence-tension-compliance}
\end{equation}

La ecuaci\'on de Einstein adopta la forma constitutiva

\begin{equation}
 G_{\mu\nu}+\Lambda g_{\mu\nu}
 =8\pi\mathscr C_P T_{\mu\nu}.
 \label{rm:eq:confluence-einstein-constitutive}
\end{equation}

Para una carga \(q\), la coordenada el\'ectrica elevada al plano de energ\'ia
es

\begin{equation}
 \mathcal Q(q)=Q_P\sqrt{\frac{Z}{4\pi\hbar}}\,q.
 \label{rm:eq:confluence-electric-coordinate}
\end{equation}

Entonces Newton y Coulomb se re\'unen en la forma bilineal

\begin{equation}
 \boxed{
 V(r)=\frac{\mathscr C_P}{r}
 \bigl(\mathcal Q_1\mathcal Q_2-E_1E_2\bigr).}
 \label{rm:eq:confluence-newton-coulomb}
\end{equation}

El cono \(\mathcal Q^2=E^2\) reproduce la condici\'on est\'atica de
extremalidad en la carta Einstein--Maxwell.  El producto de hojas del vac\'io
gobierna el cono causal; su cociente gobierna la elevaci\'on el\'ectrica y la
orientaci\'on.

La identidad cruzada con el sector electrod\'ebil es

\begin{equation}
 \boxed{
 \frac{\hbar^2}{8\pi\ell_P^2}
 \frac{(\kappa/\mu)}{e^2}
 =\frac1{4\pi\alpha_{\HMT}}
 =\frac1{g^2}+\frac1{g'^2}.}
 \label{rm:eq:confluence-einstein-maxwell-ew}
\end{equation}

Esta igualdad conserva su valor bajo la inversi\'on de la hoja de acci\'on:
\(G\) y \(e^2\) transforman contragredientemente, mientras
\(G\hbar\), \(Ge^2\) y \(\ell_P^2\) permanecen invariantes.

\subsection{3.ª confluencia: geometría, información y mezcla}

Sea \(N=N_{90}+N_{120}\) el n\'umero total de excitaciones de borde y

\begin{equation}
 D_{90/120}=90N_{90}+120N_{120}.
 \label{rm:eq:confluence-modular-provenance}
\end{equation}

El \`area lee \(N\), mientras el Hamiltoniano hologr\'afico modular lee
\(D_{90/120}\). Como

\begin{equation}
 \det\begin{pmatrix}1&1\\90&120\end{pmatrix}=30,
 \label{rm:eq:confluence-modular-determinant}
\end{equation}

la pareja reconstruye exactamente

\begin{equation}
 N_{90}=4N-\frac{D_{90/120}}{30},
 \qquad
 N_{120}=\frac{D_{90/120}}{30}-3N.
 \label{rm:eq:confluence-sector-reconstruction}
\end{equation}

La geometr\'ia total y la energ\'ia modular ponderada conservan, por tanto,
cantidad y procedencia.

El estado bicapa angular

\begin{equation}
 \rho_y=\frac{e^{-yR}}{2\cosh y}
 \label{rm:eq:confluence-rhoy}
\end{equation}

posee entrop\'ia par y divergencia orientada

\begin{equation}
 S(\rho_y)=\log(2\cosh y)-y\tanh y,
 \qquad
 D(\rho_y\Vert\rho_{-y})=2y\tanh y.
 \label{rm:eq:confluence-info-orientation}
\end{equation}

La involuci\'on \(y\mapsto-y\) se representa en el sector de sabor mediante

\begin{equation}
 R_{\rm CKM}
 =M_{\rm CKM}\operatorname{diag}(1,-1,1)M_{\rm CKM}^{-1},
 \qquad
 R_{\rm CKM}^2=I,
 \qquad
 \det R_{\rm CKM}=-1.
 \label{rm:eq:confluence-rckm}
\end{equation}

Barbero lee la componente impar de la transici\'on \(6\to8\), la entrop\'ia
lee su distribuci\'on par y \(R_{\rm CKM}\) realiza la misma inversi\'on en el
codominio de mezcla.  Son representaciones coordinadas de una bisagra com\'un,
con dominios y observables propios.

\subsection{4.ª confluencia: pantalla, torsión y memoria helicoidal}

En el cociente lorentziano \(Q_{\square}\), la soldadura y la respuesta
positiva satisfacen

\begin{equation}
 \boxed{
 (\beta_m^{\square})^*\Lambda_m^{\square}\beta_m^{\square}
 =\frac{56}{81}I_{Q_{\square}}.}
 \label{rm:eq:confluence-screen-energy}
\end{equation}

En el registro profundo, el operador de tornillo verifica

\begin{equation}
 \boxed{
 V^{-1}\mathscr SV=(2+\sqrt3)\ii\,\mathscr S.}
 \label{rm:eq:confluence-screw}
\end{equation}

La primera identidad mide el coste m\'inimo de respuesta de la pantalla; la
segunda transporta orientaci\'on y profundidad. Catalan aparece como el
segundo momento del componente cuaternario. Si
\(N_{\rm dep}\delta_n=n\delta_n\) denota el operador de profundidad del
registro helicoidal, entonces

\begin{equation}
 G_{\rm Cat}=\Tr(N_{\rm dep}^{-2}\mathcal Q_{\chi,-4}).
 \label{rm:eq:confluence-catalan}
\end{equation}

La ley de Perron

\begin{equation}
 M_n=\left(\frac{a_n}{a_0}\right)^{-6}
 \label{rm:eq:confluence-sixth-law}
\end{equation}

ofrece la covariancia exacta de una densidad cuadr\'atica de esp\'in.  El
realizador de Cartan debe transportar
\cref{rm:eq:confluence-screen-energy,rm:eq:confluence-screw} a una corriente y una
contorsi\'on cuyo modo homog\'eneo satisfaga
\cref{rm:eq:confluence-sixth-law}.  Este encadenamiento sit\'ua la interfaz
f\'isica en una flecha \`unica y falsable.

\subsection{Confluencia entre periodicidad del reloj, energía de horizonte y coste informacional}

El reloj \(108\) fija

\begin{equation}
 E_P=k_B\Theta_{\rm clk}\log3.
 \label{rm:eq:confluence-planck-trit}
\end{equation}

Para el parche de Sitter de radio \(R=\ell_P\), la energ\'ia de vac\'io, la
temperatura y la entrop\'ia cumplen

\begin{equation}
 E_\Lambda=T_H S_H=\frac{E_P}{2}
 =\frac{k_B\Theta_{\rm clk}\log3}{2}.
 \label{rm:eq:confluence-horizon-half}
\end{equation}

De aqu\'i resultan

\begin{equation}
 E_\Lambda t_P=\frac{\hbar}{2},
 \qquad
 E_\Lambda T_{108}=\frac{h}{2},
 \qquad
 \frac{S_H}{k_B}=\pi.
 \label{rm:eq:confluence-half-action}
\end{equation}

El volumen del vac\'io, el \`area del horizonte, la temperatura del ciclo y el
coste informacional de un trit publican la misma energ\'ia.  La multiplicidad
\(e^{S_H/k_B}=e^\pi\) coincide exactamente con el autovalor expansivo de
\(e^{\pi J_\varphi}\); la promoci\'on de esta igualdad a equivalencia de
microestados exige un entrelazador que conserve medida y din\'amica.

\subsection{Confluencia fibrada de las respuestas HMT–MD: vacío, acción, gravedad, información, mezcla, memoria y cosmología}

\begin{theorem}[Confluencia de respuestas HMT--MD]
Fijado el estado enriquecido \(X\) y las primitivas declaradas en cada lector,
las publicaciones de vac\'io, acci\'on, gravedad, informaci\'on, mezcla,
pantalla, memoria helicoidal y cosmolog\'ia son las proyecciones o
composiciones definidas sobre el producto fibrado de grafos
\cref{rm:eq:master-fiber-product}. Sus relaciones cruzadas quedan
generadas por
\cref{rm:eq:confluence-ew-master,rm:eq:confluence-newton-coulomb,%
rm:eq:confluence-einstein-maxwell-ew,rm:eq:confluence-sector-reconstruction,%
rm:eq:confluence-info-orientation,rm:eq:confluence-screen-energy,%
rm:eq:confluence-screw,rm:eq:confluence-horizon-half}.
Cada evaluaci\'on num\'erica es terminal y cada memoria orientada permanece en
el factor que la genera.
\end{theorem}

\begin{proof}
La construcci\'on de \(\mathfrak M_X\) fija el antecedente com\'un. Las
ecuaciones de vac\'io siguen del c\'alculo funcional de \(T\); las ecuaciones
electrod\'ebiles siguen de \(D\) y \(Q(D)\); la forma electrogravitatoria sigue
de las dos identidades de la pareja Planck; la reconstrucci\'on \(90/120\)
resulta de invertir una matriz entera de determinante \(30\); la orientaci\'on
informacional sigue del logaritmo de \(\rho_y\); la energ\'ia de pantalla
resulta de la respuesta \`unica en \(Q_{\square}\); el tornillo sigue de las
dos covariancias de Weyl y Pell; y la igualdad de horizonte sigue de las
f\'ormulas geom\'etricas de Sitter especializadas en \(R=\ell_P\). Ninguna
prueba utiliza el decimal terminal para seleccionar la rama.  El teorema de
independencia horizontal impide convertir las relaciones cruzadas en
factorizaciones no demostradas.
\end{proof}

\subsection{Corolario de confluencia metrológica entre las realizaciones HMT y SI}

Toda salida admite la qu\'intupla inseparable

\begin{equation}
 \mathcal K(P)=
 \bigl(
 \text{antecedente},
 \text{ecuaci\'on},
 \text{invariante},
 \text{significado},
 \text{valor terminal}
 \bigr).
 \label{rm:eq:confluence-quintuple}
\end{equation}

El atlas metrol\'ogico del volumen enumera estas qu\'intuplas.  Su contenido
probatorio es el siguiente: el antecedente fija la genealog\'ia; la ecuaci\'on
fija la operaci\'on; el invariante demuestra qu\'e sobrevive; el significado
explicita la lectura f\'isica; el valor terminal permite el reconocimiento
externo.  Separar una columna destruye informaci\'on causal; conservar las
cinco permite insertar el resultado en una carta de unidades sin convertir
esa carta en fundamento.

\begin{proofstatusbox}
El producto fibrado, las identidades algebraicas y el teorema de independencia
horizontal constituyen una formalizaci\'on exacta.  Las ecuaciones internas
citadas conservan el estatuto demostrado en los cap\'itulos que contienen sus demostraciones completas.  Las
lecturas Einstein--Maxwell, de Sitter, Einstein--Cartan y electrod\'ebil se
mantienen en sus interfaces declaradas cuando requieren un realizador
adicional.  El teorema organiza esas capas sin promoverlas por similitud
num\'erica.
\end{proofstatusbox}

\begin{falsifierbox}
La confluencia se refuta si dos componentes atribuidas al mismo elemento de
\(\mathfrak M_X\) proceden realmente de estados incompatibles, si falla una
de las identidades generadoras, si una evaluaci\'on terminal interviene en la
selecci\'on de su propio antecedente, o si se pierde una memoria que otro
lector distingue.  Una pretendida factorizaci\'on horizontal queda refutada
por cualquier par de estados con igual lector intermedio y salida final
distinta.  Las promociones f\'isicas se contrastan, adem\'as, con sus propios
realizadores, identidades de gauge, Bianchi, refinamiento y condiciones de
frontera.
\end{falsifierbox}
